% 4.4 O4 Synthesis (all). Source: deck slide 33 (evidence-coded matrix), with Q/D/N codes kept as presented.
\subsection{O4: Synthesis by boundary and scale}\label{sec:o4}

\begin{table*}[t]\centering
\caption{Each pressure has its own scale. Planetary boundaries against the scale of impact of AI data centres.
\textbf{Q}: quantified in this project; D: documented pathway, not quantified; N: not assessed. Status is the 2025
global assessment of all human activity \citep{sakschewski2025}, not a finding that data centres caused it.}
\label{tab:matrix}
\footnotesize\renewcommand{\arraystretch}{1.12}
\begin{tabular}{@{}P{2.35cm}P{1.7cm}P{3.75cm}P{3.75cm}P{3.75cm}@{}}
\toprule
Boundary & Status (2025) & Local & Regional & Global \\
\midrule
Climate change & Transgressed & D\enspace Backup generators; construction & \textbf{Q}\enspace Grid intensity 212--559\,g
\cotwo/kWh (\Fig{fig:o1}) & D\enspace Operational and embodied greenhouse gases \\
Freshwater change & Transgressed & \textbf{Q}\enspace 34\% of capacity in stressed basins; Mesa up to 15\,M\,L/day
& \textbf{Q}\enspace 67--86\% of water used at power plants & \textbf{Q}\enspace $<$0.1\% of Earth3's limit: the
total hides the hotspots \\
Land-system change & Transgressed & D\enspace Campus and substation footprint & D\enspace Transmission corridors;
mining land & D\enspace Land conversion in the supply chain \\
Biosphere integrity & Transgressed & D\enspace Habitat effects of sites & D\enspace Mining, river and wetland pressure &
D\enspace Cumulative ecosystem pressure \\
Biogeochemical flows & Transgressed & N & D\enspace Generation, mining and runoff & N \\
Novel entities & Transgressed & D\enspace Refrigerants, chemicals, exhaust & \textbf{Q}\enspace E-waste collected:
42.8\% in Europe, 0.7\% in Africa & D\enspace Hazardous material flows \\
Ocean acidification & Transgressed & N & N & D\enspace Through \cotwo\ only \\
Aerosol loading & Safe & D\enspace Diesel and construction particulates & D\enspace Combustion pollution & N \\
Stratospheric ozone & Safe & N & N & N \\
\bottomrule
\end{tabular}
\end{table*}

\Tab{tab:matrix} brings the four objectives together. The strongest evidence is for freshwater, which we quantified
at all three scales: locally (a third of capacity in stressed basins, and Mesa's contract), regionally (most of the
water is evaporated at power plants on the regional grid) and globally (a negligible share of the planetary limit,
which is itself the finding). Climate change acts through the grid: the load is regional and the carbon intensity is set
regionally, but the emissions add to a global stock. Minerals and e-waste move through global supply chains, so their
impacts are displaced: mined in one region, used in a second and discarded in a third. Most other cells are documented
pathways that we did not measure; we mark them as such rather than fill them. The siting audit adds who decides:
at Visakhapatnam no instrument in force required a water test before the site was chosen (Section~\ref{sec:siting}), so the
local water pressure shown in the matrix was not tested at the site before approval.

\begin{reading}
Our reading: the matrix also shows why a safe-and-just view \citep{rockstrom2023} matters for this topic. A boundary
defined only by the global total says AI's water is safe. A boundary that asks who is harmed points to the 21 million
people near stressed sites, a desert city on a shrinking river, and a port city already short of drinking water.
\end{reading}

\begin{finding}{Finding 6. Small globally, concentrated locally}
AI's footprint is a small share of the world's electricity and water, but it piles up in a few countries, grids,
basins and towns. Each pressure has its own scale, and each scale needs its own response.
\end{finding}
